\section{Research context and objective}
Drug--target affinity prediction estimates the strength of association between a molecule and a protein. Binding affinity is continuous, often reported as a dissociation constant $K_d$; lower $K_d$ corresponds to stronger binding. Predicting a numeric score preserves more information than assigning a binary interacting/non-interacting label, though a prediction remains a computational estimate rather than an assay result \citep{ozturk2018deepdta,yang2022mgraphdta}. The project objective is to implement and evaluate a graph-plus-sequence MGraphDTA-style model on Davis, make the data and optimization pipeline inspectable, and demonstrate attribution and structure-analysis utilities.

The original MGraphDTA work introduced dense multiscale graph encoding for molecular structure, multiscale sequence convolutions for targets, and Grad-AAM for atom-level visualization. It reports Davis MSE $0.265\pm0.002$, CI $0.822\pm0.001$ and $r_m^2=0.701\pm0.001$ for its reported benchmark setting \citep{yang2022mgraphdta}. These are paper-reported values, not results from this project. Earlier sequence-based work such as DeepDTA established a continuous affinity-prediction formulation using learned sequence representations \citep{ozturk2018deepdta}. The present implementation is compared with the method and author code \citep{mgraphdtaauthorscode}; it is not assumed to reproduce their protocol.

\section{Data and preprocessing}
The local Davis sources are compound and target dictionaries, a pickled Kd matrix, and the supplied fold-index files under \texttt{data/raw/davis/}. This agrees with the Davis kinase-selectivity study and the benchmark use described in DeepDTA \citep{davis2011,davisdeepdta,ozturk2018deepdta}. The project builds one row per compound--target pair. Its verified counts are 68 compounds, 442 target IDs, 379 unique protein sequences and 30,056 records. The official train indices flatten to 25,046 interactions; the independent test indices contain 5,010. Training further selects a seeded 10\% validation subset from the training fold: 2,505 validation and 22,541 optimization records in the primary setup. This is an interaction-level split, not a drug-disjoint or target-disjoint evaluation.

The label transform is
\begin{equation}\label{eq:pkd}pK_d=-\log_{10}(K_d/10^9)=9-\log_{10}(K_d\text{ in nM}).\end{equation}
The Davis convention maps 10,000 nM labels to p$K_d=5$. No additional target normalization is applied in the default path. MSE is therefore measured in squared p$K_d$ units and RMSE in p$K_d$ units.

RDKit parses SMILES and produces one node per atom. Each atom has 22 features: 9 element indicators (H, C, N, O, F, Cl, S, Br, I); atomic number; RDKit acceptor and donor flags; aromaticity; 3 hybridization indicators (SP, SP2, SP3); total hydrogen count; explicit valence; formal charge; implicit valence; explicit hydrogen count; and radical electron count. Each undirected bond becomes two directed entries in `edge_index`. No bond feature tensor is sent to this project’s model. Invalid/empty SMILES raise an error rather than being silently removed.

Protein sequences use a 25-symbol vocabulary with zero for padding and unknown symbols. Input is uppercased, padded to 1,200 residues or truncated at 1,200; the longest Davis sequence is 2,549 residues. Unknown symbols and padding share ID zero. The model does not pass an explicit convolution mask, and the preprocessing code does not filter nonfinite/nonpositive Kd values before the logarithm; these are limitations of the implemented handling.

\section{Model implementation}
\begin{figure}[tb]\centering
\begin{tikzpicture}[node distance=5mm and 9mm,box/.style={draw=navy,rounded corners,fill=light,align=center,text width=31mm,minimum height=9mm,font=\small},arr/.style={-{Latex[length=2mm]},thick,draw=navy}]
\node[box](sm){Drug SMILES};\node[box,below=of sm](graph){RDKit graph\\22-D atom nodes + edges};\node[box,below=of graph](mgnn){MGNN\\3 dense blocks, 8 units each};\node[box,below=of mgnn](drug){Global mean pool\\96-D drug vector};
\node[box,right=34mm of sm](seq){Protein sequence};\node[box,below=of seq](tok){Tokenize\\pad/truncate to 1,200};\node[box,below=of tok](mcnn){Embedding + MCNN\\3 scales (RF 3, 5, 7)};\node[box,below=of mcnn](protein){Max pool + projection\\96-D protein vector};
\node[box,below=16mm of drug,xshift=17mm](fusion){Concatenate\\192 dimensions};\node[box,below=of fusion](head){MLP\\1024 $\to$ 1024 $\to$ 256};\node[box,below=of head](out){Predicted p$K_d$};
\foreach \a/\b in {sm/graph,graph/mgnn,mgnn/drug,seq/tok,tok/mcnn,mcnn/protein,fusion/head,head/out}{\draw[arr](\a)--(\b);}
\draw[arr](drug.south)--++(0,-5mm)-|(fusion.west);\draw[arr](protein.south)--++(0,-5mm)-|(fusion.east);
\end{tikzpicture}
\caption{Main forward path as implemented in this repository. The protein input is a sequence, not a 3D structure.}
\label{fig:architecture}\end{figure}

\paragraph{MGNN.} A GraphConv stem maps 22 atom channels to 32, followed by node-level batch normalization and ReLU. Each of three dense blocks has eight units; every unit applies GraphConv--BN--ReLU to a bottleneck width of 64 and then to 32 growth channels, concatenating the new feature map with prior features. A transition layer halves the channels after each block. The progression is 32 $\to$ 288 $\to$ 144 $\to$ 400 $\to$ 200 $\to$ 456 $\to$ 228. Global mean pooling over atoms followed by a linear map yields a 96-D molecule vector. The code’s “27 depth” convention counts 24 dense units and three transitions, excluding conv0; it has 28 top-level stages including conv0 and 52 literal GraphConv modules.

\paragraph{MCNN and regression.} A trainable embedding maps token IDs to 128 dimensions. Three parallel branches apply one, two or three Conv1D layers (96 channels, width 3, stride 1, no padding, ReLU), giving receptive fields 3, 5 and 7. Adaptive max pooling produces 96 values per branch; concatenation and a linear layer project 288 values to 96. Drug and protein vectors concatenate to 192 values and enter a head of linear widths 1024, 1024, 256 and 1, with ReLU and dropout 0.1 after each of the first three layers. The final scalar is predicted p$K_d$. Training uses MSE regression.

\paragraph{Paper/release differences.} The author release has a closely aligned default MGNN/MCNN/head topology. However, its regression preprocessing min--max scales each molecule’s whole atom-feature matrix using one scalar minimum and maximum. This project’s default is unscaled; an optional versioned scaled bundle and separate scaled run exist. Training protocols also differ: the primary run used batch size 32 and a 150-epoch cap with manual interruption; the paper configuration is documented as batch size 512 and learning rate 0.0005. The local paper PDF exists, but text extraction was blocked by local MiKTeX first-run setup, and full ESI verification remains incomplete. Accordingly, architecture correspondence is based on source code and local compliance notes, not a claim of complete replication.

\section{Training and evaluation}
The main trainer defaults are batch size 512, learning rate 0.0005, 1,000 epochs, 50 steps per epoch, patience 400, validation fraction 0.1, seed 42. Run-specific metadata overrides defaults. The trainer uses Adam without explicit weight decay, `nn.MSELoss`, no scheduler and no gradient clipping in the main fit path. Best checkpoint is selected on validation MSE; test data is evaluated after fit. Device is recorded as CUDA for the principal runs. The checked environment documented in the project plan is Python 3.12, PyTorch 2.6.0+cu124, PyG 2.8.0 and RDKit 2026.03.6; requirements pins PyG 2.8.0.post1. These are documented versions, not a live capture from this audit.

MSE and RMSE measure absolute error. CI measures pairwise ordering agreement, giving prediction ties half credit and excluding true-target ties. It does not measure calibration. The implementation’s modified squared correlation uses Pearson $r^2$, $k=\sum y\hat y/\sum\hat y^2$, zero-intercept $r_0^2$, and $r_m^2=r^2[1-\sqrt{|(r^2)^2-(r_0^2)^2|}]$ \citep{roy2009}. This is not ordinary $R^2$. The earlier primary `evaluation_metrics.json` contains $r_m^2=0.3037$; the corrected file and current metric implementation give 0.6817, which is used below.

